\lecture{37}{Непрерывная распределённая оптимизация}

\sourcevideo
    {https://www.youtube.com/playlist?list=PLOzRYVm0a65fhKDS8cYIC7YCuVGJ4FOJx}
    {Week 10: Lecture 37}

Рассматривается задача
\[
    \min_{x\in\mathbb R^d}F(x),
    \qquad
    F(x)=\sum_{i=1}^{N}f_i(x),
\]
в которой агент \(i\) знает только функцию \(f_i\) и обменивается
сообщениями с соседями. Цель --- построить непрерывную систему,
достигающую единственного минимума за время, равномерно ограниченное
по всем начальным состояниям.

\section{Централизованный поток}

Предположим, что функции \(f_i\) дифференцируемы и выпуклы, а сумма
\(F\) является \(\mu\)-сильно выпуклой. Тогда минимум \(x^\star\)
единственен и
\[
    \lVert\nabla F(x)\rVert^2
    \ge
    2\mu\bigl(F(x)-F(x^\star)\bigr).
\]
Для \(z\in\mathbb R^d\) и \(\nu>0\) положим
\[
    \operatorname{sig}^{\nu}(z)
    =
    \begin{cases}
        z\lVert z\rVert^{\nu-1},&z\ne0,\\[1mm]
        0,&z=0.
    \end{cases}
\]
Тогда
\[
    z^{\mathsf T}\operatorname{sig}^{\nu}(z)
    =
    \lVert z\rVert^{1+\nu}.
\]
Выберем
\[
    a_1,a_2>0,
    \qquad
    0<\nu_1<1<\nu_2,
\]
и определим
\[
    \Psi(z)
    =
    a_1\operatorname{sig}^{\nu_1}(z)
    +
    a_2\operatorname{sig}^{\nu_2}(z).
\]
Централизованный поток имеет вид
\[
    \dot x=-\Psi\bigl(\nabla F(x)\bigr).
\]
Для
\[
    V(x)=F(x)-F(x^\star)
\]
получаем
\[
\begin{aligned}
    \dot V
    ={}&
    -a_1\lVert\nabla F(x)\rVert^{1+\nu_1}
    \\
    &-
    a_2\lVert\nabla F(x)\rVert^{1+\nu_2}.
\end{aligned}
\]
Положим
\[
    p_1=\frac{1+\nu_1}{2},
    \qquad
    p_2=\frac{1+\nu_2}{2}.
\]
Тогда \(0<p_1<1<p_2\) и
\[
    \dot V
    \le
    -c_1V^{p_1}-c_2V^{p_2},
\]
где
\[
    c_1=a_1(2\mu)^{p_1},
    \qquad
    c_2=a_2(2\mu)^{p_2}.
\]
Следовательно,
\[
    T_F
    \le
    \frac1{c_1(1-p_1)}
    +
    \frac1{c_2(p_2-1)}.
\]
Эта граница не зависит от начальной точки.

\section{Недоступность глобального градиента}

В распределённой сети агенты начинают из произвольных точек
\[
    x_i(0)\ne x_j(0)
\]
и агент \(i\) может вычислить только \(\nabla f_i(x_i)\), но не
\[
    \nabla F(x)=\sum_{j=1}^{N}\nabla f_j(x).
\]
Поэтому требуется одновременно согласовать состояния \(x_i\) и
восстановить текущий средний локальный градиент
\[
    \theta_{\mathrm c}(t)
    =
    \frac1N
    \sum_{i=1}^{N}\nabla f_i(x_i(t)).
\]
Каждый агент хранит оценку \(\theta_i(t)\in\mathbb R^d\). Если
локальная функция дважды дифференцируема, производная её градиента
равна
\[
    h_i(t)
    =
    \frac{\mathrm d}{\mathrm dt}\nabla f_i(x_i(t))
    =
    \nabla^2f_i(x_i(t))\dot x_i(t).
\]
Таким образом, точный непрерывный трекинг требует вычисления
произведения гессиана на вектор либо иного способа получить
производную локального градиента.

\section{Фиксированно-временной трекинг градиента}

Инициализация задаётся условием
\[
    \theta_i(0)=\nabla f_i(x_i(0)).
\]
Далее
\[
    \dot\theta_i=h_i+\omega_i,
\]
где для симметричных весов \(a_{ij}=a_{ji}>0\)
\[
\begin{aligned}
    \omega_i
    ={}&
    -k_0\sum_{j\in\mathcal N_i}
    a_{ij}\operatorname{sign}(\theta_i-\theta_j)
    \\
    &-
    k_1\sum_{j\in\mathcal N_i}
    a_{ij}\operatorname{sig}^{\gamma_1}(\theta_i-\theta_j)
    \\
    &-
    k_2\sum_{j\in\mathcal N_i}
    a_{ij}\operatorname{sig}^{\gamma_2}(\theta_i-\theta_j),
\end{aligned}
\]
с параметрами
\[
    k_0,k_1,k_2>0,
    \qquad
    0<\gamma_1<1<\gamma_2.
\]
Нечётность парных взаимодействий даёт
\[
    \sum_{i=1}^{N}\omega_i=0.
\]
Следовательно,
\[
\begin{aligned}
    \frac{\mathrm d}{\mathrm dt}
    \sum_{i=1}^{N}\theta_i
    &=
    \sum_{i=1}^{N}h_i
    \\
    &=
    \frac{\mathrm d}{\mathrm dt}
    \sum_{i=1}^{N}\nabla f_i(x_i).
\end{aligned}
\]
С учётом инициализации сохраняется инвариант
\[
    \sum_{i=1}^{N}\theta_i(t)
    =
    \sum_{i=1}^{N}\nabla f_i(x_i(t)).
\]
Поэтому после достижения консенсуса оценки обязаны совпасть с
\(\theta_{\mathrm c}(t)\).

Пусть для всех рёбер и всех \(t\)
\[
    \lVert h_i(t)-h_j(t)\rVert
    \le
    \rho_\theta.
\]
При достаточно большом \(k_0\), зависящем от этой границы и графа,
знаковый член подавляет возмущение, а два степенных члена обеспечивают
существование времени \(T_\theta\), не зависящего от начальных
оценок, такого что
\[
    \theta_i(t)=\theta_{\mathrm c}(t),
    \qquad
    t\ge T_\theta.
\]
Порог \(k_0\) выбирается так, чтобы используемые далее оценки выполнялись для всех \(k\ge k_0\).

\section{Согласование локальных переменных}

Используя оценку среднего градиента, агент строит оптимизационное
направление
\[
    g_i=-\Psi(N\theta_i).
\]
Для согласования переменных добавляется сигнал
\[
\begin{aligned}
    u_i
    ={}&
    -\ell_0\sum_{j\in\mathcal N_i}
    a_{ij}\operatorname{sign}(x_i-x_j)
    \\
    &-
    \ell_1\sum_{j\in\mathcal N_i}
    a_{ij}\operatorname{sig}^{\delta_1}(x_i-x_j)
    \\
    &-
    \ell_2\sum_{j\in\mathcal N_i}
    a_{ij}\operatorname{sig}^{\delta_2}(x_i-x_j),
\end{aligned}
\]
где
\[
    \ell_0,\ell_1,\ell_2>0,
    \qquad
    0<\delta_1<1<\delta_2.
\]
Полная динамика равна
\[
    \dot x_i=g_i+u_i.
\]
До завершения трекинга направления \(g_i\) могут различаться и
действуют как возмущение консенсусной подсистемы. Предположим, что
\[
    \lVert g_i(t)-g_j(t)\rVert\le\rho_x
\]
на каждом ребре. Если \(\ell_0\) достаточно велик относительно
\(\rho_x\) и параметров графа, то существует равномерное время
\(T_x\), после которого
\[
    x_1(t)=\cdots=x_N(t).
\]
Знаковый член также обеспечивает инвариантность консенсусного
множества при продолжающемся воздействии различающихся направлений.

\section{Сведение к централизованному потоку}

Положим
\[
    T_0=\max\{T_\theta,T_x\}.
\]
При \(t\ge T_0\) одновременно выполнены
\[
    \theta_i
    =
    \frac1N\sum_{j=1}^{N}\nabla f_j(x_j)
\]
и
\[
    x_1=\cdots=x_N=x.
\]
Поэтому
\[
    N\theta_i
    =
    \sum_{j=1}^{N}\nabla f_j(x)
    =
    \nabla F(x),
\]
а консенсусный сигнал удовлетворяет \(u_i=0\). Все агенты следуют
одному уравнению
\[
    \dot x=-\Psi\bigl(\nabla F(x)\bigr).
\]
После момента \(T_0\) до достижения \(x^\star\) требуется не более
\(T_F\). Следовательно,
\[
    T_{\mathrm{total}}
    \le
    \max\{T_\theta,T_x\}+T_F.
\]

\begin{theorem}[Фиксированно-временная распределённая оптимизация]
Пусть граф связен, неориентирован и имеет симметричные положительные
веса. Пусть \(F=\sum_i f_i\) сильно выпукла, локальные функции
достаточно гладки, различия \(h_i-h_j\) и \(g_i-g_j\) равномерно
ограничены, а робастные коэффициенты выбраны выше соответствующих
порогов. Тогда существует число \(T_{\max}<\infty\), не зависящее от
начальных состояний, такое что
\[
    x_i(t)=x^\star
\]
для всех агентов и всех \(t\ge T_{\max}\).
\end{theorem}

\section{Реализация и границы модели}

Каждый агент передаёт соседям два вектора размерности \(d\):
\(x_i(t)\) и \(\theta_i(t)\). Поэтому объём одного полного обмена
имеет порядок
\[
    O(|\mathcal E|d),
\]
но с примерно вдвое большей постоянной, чем при передаче только
состояний. Локально требуются \(\nabla f_i(x_i)\) и
\[
    h_i=\nabla^2f_i(x_i)\dot x_i.
\]

Функция \(\operatorname{sign}\) разрывна в нуле, поэтому строгий
анализ ведётся в смысле решений Филиппова. Гладкое насыщение уменьшает
дребезг, но обычно заменяет точный консенсус попаданием в малую
окрестность.

Результат предполагает непрерывное время, статический связный
неориентированный граф, точные сообщения без задержек, доступность
производной локального градиента и известные границы возмущений. Для
ориентированных, асинхронных и изменяющихся сетей требуется другая
схема трекинга. Прямая дискретизация также не обязана сохранять
точную фиксированно-временную гарантию.
